% FILE: 00_project_identity.tex
% Project: otel-weaver-exp-002-diff-to-residual
% Thesis Role: weaver-diff-to-pi-residual-bridge-experiment

\section{Project Identity}
\label{sec:otel-weaver-exp-002-diff-to-residual-identity}

\subsection{Purpose}
\label{sec:otel-weaver-exp-002-diff-to-residual-identity-purpose}

Experiment EXP-002 establishes a formal, executable boundary between design-time
telemetry schema evolution and runtime process behavioral deviation. Its immediate
engineering purpose is to manufacture the Rust library crate
\texttt{exp-002-weaver-diff-to-pi-residual}, which implements a schema-translation
bridge called \texttt{BridgeRx}. The bridge intercepts incoming telemetry feedstock
that conforms to schema version S2 and translates it into the S1-compatible structure
that the conformance court's validators expect, before any conformance algorithm fires.

The broader scientific purpose is to prove that the PI Conformance Residual
$\mathcal{R}_{\Delta} = |f_{S_1}(\mathcal{L}_2, M) - f_{S_2}(\mathcal{L}_2, M)|$
can be driven to zero by a lawfully placed translation layer. Without this bridge, a
schema attribute rename---such as the Petri-net-motivated change from
\texttt{process.pi.activity.name} to \texttt{process.pi.transition.name}---would
cause the conformance court to misclassify a structurally clean telemetry record as a
process violation, triggering a false compliance alert. EXP-002 proves that this
category of false positive is preventable at the architectural boundary rather than
requiring post-hoc court correction.

The experiment is the second of five in the OTel Weaver series
(exp-001 through exp-005). Its upstream dependency is the custom process intelligence
semantic convention registry defined in exp-001. Its downstream consumers include
exp-003 (live-check to refusal) and exp-004 (registry-to-witness code generation),
which both assume that telemetry entering the conformance boundary is already
normalized to the court's expected schema surface.

\subsection{Repository Surface}
\label{sec:otel-weaver-exp-002-diff-to-residual-identity-repository}

The experiment lives at:
\begin{verbatim}
/Users/sac/process-intelligence/otel-weaver/experiments/
    exp-002-weaver-diff-to-pi-residual/
\end{verbatim}

The Rust library crate is declared in
\texttt{Cargo.toml} with package name \texttt{exp-002-weaver-diff-to-pi-residual},
library target \texttt{bridge}, and source path \texttt{src/bridge.rs}. Declared
dependencies are \texttt{serde} (v1.0, derive feature), \texttt{serde\_json} (v1.0),
and \texttt{serde\_yaml} (v0.9). The Cargo dependency graph was materialized to
\texttt{target/debug/deps} including compiled \texttt{serde}, \texttt{serde\_json},
\texttt{serde\_yaml}, \texttt{indexmap}, and \texttt{hashbrown} rlibs, confirming that
\texttt{cargo build} was executed at some point in the experiment's history.

Parent-scope checkpoints are held at
\texttt{/Users/sac/process-intelligence/otel-weaver/checkpoints/} and cover the full
five-experiment suite rather than individual experiments. No receipt files are present
inside the exp-002 directory itself; the checkpoint scope is intentionally aggregated
at the \texttt{otel-weaver/} level.

\subsection{Primary Research Role}
\label{sec:otel-weaver-exp-002-diff-to-residual-identity-role}

Within the dissertation, EXP-002 occupies the role of
\emph{weaver-diff-to-pi-residual-bridge-experiment}. Its specific scientific
contribution is the formal categorical separation encoded in the invariant:
\[
    \Delta_W(S_1, S_2) \neq \delta_P(\mathcal{L}_1, \mathcal{L}_2)
\]
This law states that a Weaver Diff---a set of \texttt{(path, op, value)} tuples where
$op \in \{\text{Add}, \text{Remove}, \text{Rename}, \text{TypeChange}\}$---belongs to
the schema metadata plane and must never be evaluated inside the process behavioral
deviation court. The experiment manufactures an executable proof of this law as a
Rust type system and runtime translation bridge rather than as an informal assertion.

EXP-002 directly supports the thesis's claim that full-lifecycle process intelligence
requires strict nominal category separation enforced at every ingestion boundary, and
that engineering mechanisms (not documentation conventions) must maintain those
categories under schema evolution pressure.

\subsection{Key Artifacts}
\label{sec:otel-weaver-exp-002-diff-to-residual-identity-artifacts}

The following artifacts constitute the experiment's evidence surface, as extracted
from the repository:

\begin{itemize}
    \item \texttt{src/bridge.rs} --- Single source file implementing \texttt{BridgeRx},
          \texttt{TelemetryFeedstockS1}, \texttt{TelemetryFeedstockS2}, and the inline
          unit test \texttt{test\_bridge\_translation\_equivalence}.
    \item \texttt{Cargo.toml} --- Library crate manifest declaring the \texttt{bridge}
          target and three serde-family dependencies.
    \item \texttt{README.md} --- Mathematical framework document defining
          $\Delta_W$, $\delta_P$, and $\mathcal{R}_{\Delta}$ with worked JSON diff
          examples and Rust code listings.
    \item \texttt{target/debug/deps/} --- Materialized dependency artifacts confirming
          \texttt{cargo build} was executed.
    \item \texttt{../checkpoints/GGEN\_OTEL\_WEAVER\_PI\_ALIVE\_001.md} --- Parent-scope
          checkpoint certifying the full five-experiment suite as COMPLETE \& ACTIVE
          (SHA-256: \texttt{e6bc59b8d210e309228d4bf09c0fa4db55fe310efc35ad1a89b3f021e85a67f1}).
    \item \texttt{../checkpoints/GGEN\_OTEL\_WEAVER\_PI\_RUNTIME\_001.md} --- Runtime
          integration checkpoint citing 62 passing tests across \texttt{wasm4pm},
          \texttt{wasm4pm-compat}, and \texttt{blue\_river\_dam}
          (SHA-256: \texttt{86ef45b3c63ae38007701169bcf0abbda2694a43160f3f8d3695c6255b02cb44}).
    \item \texttt{../doctrine/registry-diff-is-not-process-drift.md} --- Doctrine file
          anchored at \texttt{REGISTRY\_DIFF\_NOT\_DRIFT\_001} formalizing the categorical
          separation this experiment enforces.
    \item \texttt{../mappings/registry-diff-to-pi-drift-map.yaml} --- YAML boundary rules
          (\texttt{drift\_isolation\_rule\_1}, \texttt{drift\_isolation\_rule\_2}) routing
          schema-diff events away from the process drift court.
\end{itemize}
